\subsection{Multi-Target Oracle Construction}

The first component of EGAAT is a phase-inversion oracle $\mathcal{O}$ that
simultaneously marks an arbitrary set of $k$ target states. For a database of
size $N = 2^n$, the oracle acts as:
\begin{equation}
  \mathcal{O}\ket{x} = \begin{cases} -\ket{x} & \text{if } x \in T \\ \ket{x} & \text{otherwise} \end{cases}
\end{equation}
where $T = \{t_1, t_2, \ldots, t_k\}$ is the set of target indices. The oracle
is constructed as a sum of phase-kickback circuits, one per target, using
controlled-phase operations. Circuit depth scales as $\mathcal{O}(n \log k)$,
significantly shallower than the naive $\mathcal{O}(nk)$ construction.

\subsection{Dynamic Iteration Control via Quantum Amplitude Estimation}

Prior to the main search, EGAAT embeds a Maximum Likelihood Amplitude Estimation
(MLAE)~\cite{Suzuki2020} subroutine to estimate $k$ without classical enumeration.
MLAE runs Grover circuits for $m \in \{0, 1, 2, 4, 8, \ldots\}$ iterations and
fits the measurement statistics to the theoretical success probability:
\begin{equation}
  P_m(\theta) = \sin^2\!\left((2m+1)\theta\right), \quad \theta = \arcsin\!\sqrt{k/N}
\end{equation}
The MLE estimator $\hat{\theta}$ is found by maximising:
\begin{equation}
  \hat{\theta} = \arg\max_\theta \sum_m \left[ h_m \ln P_m(\theta) + (M-h_m)\ln(1-P_m(\theta)) \right]
\end{equation}
where $h_m$ is the hit count at iteration $m$ and $M$ is the shots per point.
Given $\hat{k} = N\sin^2(\hat{\theta})$, the optimal iteration count is:
\begin{equation}
  m^* = \left\lfloor \frac{\pi}{4} \sqrt{\frac{N}{\hat{k}}} \right\rfloor
\end{equation}

\subsection{Adaptive Diffusion Scheduling}

The standard Grover diffusion $D = 2\ket{s}\bra{s} - I$ is replaced by a
parameterised operator $D(\alpha)$:
\begin{equation}
  D(\alpha) = (1 - 2\alpha)I + 2\alpha N \ket{s}\bra{s}
\end{equation}
where $\alpha \in [0,1]$ is the amplification coefficient. The scheduling
function maps target density $\rho = \hat{k}/N$ to $\alpha$:
\begin{equation}
  \alpha(\rho) = \begin{cases}
    1.00 & \rho < 0.05 \\
    0.85 & 0.05 \leq \rho < 0.20 \\
    0.65 & 0.20 \leq \rho < 0.40 \\
    0.45 & 0.40 \leq \rho < 0.60 \\
    0.25 & \rho \geq 0.60
  \end{cases}
\end{equation}
This prevents over-rotation in dense target regimes while maintaining
maximum amplification power for sparse searches.
